\section{Introduction}

Agentic benchmarks increasingly decide which models teams deploy, yet many of them
rest on a single weak load-bearing element: an LLM judge reading a transcript and
assigning a score. A single judge introduces systematic bias, prefers verbose
output, and can favor text from its own model family. When the judge's provider
also has a model under test, the conflict is structural. The wider problem is
sharper still. The Berkeley RDI audit pointed one trivial agent at eight published
agentic benchmarks and drove seven of them to near-perfect scores without solving
any task, usually because the scoring rewarded plausible-looking output rather than
verified work~\cite{rdi2026}. A benchmark that can be gamed this way is not
measuring agent capability.

The most credible agentic benchmarks answer this by grading the world, not the
prose. tau-bench verifies a customer-service rollout by comparing the final
database state against an annotated goal state, with a binary reward and no LLM
judge in the outcome check~\cite{taubench2024}. The Berkeley Function Calling
Leaderboard (BFCL) evaluates multi-turn tool use through state-based and
response-based checks that eliminate the LLM evaluator
entirely~\cite{bfcl2024}. Terminal-Bench verifies each task with per-task tests
rather than a judge~\cite{terminalbench2026}, and SWE-bench grades a patch by
whether a repository's hidden tests pass~\cite{swebench2023}. The lesson is
consistent: where ground truth exists, check it deterministically; reserve
judgment for what genuinely needs it.

But pure deterministic checking under-measures exactly the dimensions that matter
in conversational domains. BFCL's state checks catch write and delete correctness
but not whether an agent handled an emotionally charged escalation with
appropriate empathy, asked the right clarifying question, or correctly refused an
out-of-scope request~\cite{bfcl2024}. Customer-success work lives precisely in that
soft-judgment space. The honest design is therefore hybrid: deterministic state
verification where a canonical world exists, and a judge panel for the soft
dimensions, with each clearly labeled as to which is which.

COT Bench is that hybrid. It evaluates multi-turn, tool-calling agents across two
production domains, banking and customer success, each with five scenario
categories that probe distinct failure modes. The contributions are:

\begin{enumerate}[leftmargin=*]
  \item \textbf{Hybrid, partly judge-independent efficacy.} A full third of the
  efficacy score is a deterministic check against a per-scenario ground-truth
  world: did the transfer actually land, was the fraud case filed, did an
  out-of-scope request leave the world unchanged. The remaining two thirds are two
  LLM-judge dimensions. The check needs no model and is free to compute
  (Section~\ref{sec:scoring}).

  \item \textbf{A cross-lab judge panel with published per-judge accountability.}
  Three judges from three labs, none under test, with median consensus and
  Krippendorff's $\alpha$ for reliability. Every judge's score is published, and a
  per-judge-versus-consensus delta makes any judge's favoritism, including the one
  remaining same-lab pairing, a measured quantity (Section~\ref{sec:scoring}).

  \item \textbf{Statistics a serious leaderboard owes its readers.} Scenario-level
  paired bootstrap confidence intervals on every headline number, rank bands over
  overlapping intervals, a publish-minimum scenario count, $\text{pass}^k$
  reliability, and a length-bias regression (Section~\ref{sec:stats}).

  \item \textbf{Governance designed against the two failure modes a
  single-maintainer benchmark is most exposed to:} silent judge drift and run
  cherry-picking. Runs are pre-registered before the first model call, judges are
  pinned with the actually-served model recorded, runs are never silently
  retracted, and a hash-pinned private holdout gives a public-versus-holdout
  overfitting tripwire (Section~\ref{sec:governance}).

  \item \textbf{A standing anti-gaming probe.} A deterministic do-nothing agent is
  run on demand and is expected to score near zero on both halves of efficacy;
  publishing that floor is direct evidence the scores reflect verified work, not
  the appearance of it (Section~\ref{sec:scoring}).
\end{enumerate}

\paragraph{Positioning.} COT Bench is not a re-implementation of tau-bench,
Terminal-Bench, or BFCL. It differs on three axes. First, it is explicitly hybrid:
deterministic state checks coexist with a judge panel for the soft dimensions,
rather than choosing one (tau-bench and BFCL are deterministic-only;
single-judge benchmarks are judge-only). Second, it scores production trade-offs
directly through the CLEAR composite (efficacy, cost, reliability, latency) rather
than a single success rate. Third, its judge methodology is itself a published
object of study: the per-judge deltas, the length-bias regression, and the
calibration protocol make the evaluation's own biases measurable. The benchmark
reports $\text{pass}^k$ from tau-bench, an SWE-bench-Pro-style private holdout, and
a do-nothing probe from the RDI audit, assembling lessons from across the field
rather than competing on a single one.
